Mô hình dựa trên giả thuyết rằng các nhà đầu tư coi tỷ suất lợi nhuận kỳ vọng là điều mong muốn và lo ngại rủi ro. Quá trình lựa chọn danh mục đầu tư gồm hai bước: đầu tiên là xác định các danh mục đầu tư tối ưu nằm trên đường biên hiệu quả dựa trên sự cân bằng giữa rủi ro và lợi nhuận, sau đó chọn danh mục phù hợp nhất với hàm thỏa dụng của từng nhà đầu tư tùy thuộc vào mức độ chấp nhận rủi ro và vốn đầu tư. Markowitz đã thiết lập nền tảng cho việc lựa chọn danh mục đầu tư bằng cách ánh xạ mong muốn đầu tư lên biểu đồ rủi ro - lợi nhuận. Trong biểu đồ này thì lợi nhuận được đo bằng giá trị trung bình và rủi ro được đo bằng phương sai hoặc độ lệch chuẩn của phân phối lợi nhuận. Công thức tính phương sai danh mục được xác định cụ thể như sau:
$$\sigma_w^2 = \text{Var}(R_w) = \sum_{i=1}^{N} \sum_{j=1}^{N} w_i w_j \sigma_i \sigma_j \rho_{ij}$$
Trong đó, $\rho_{ij}$ là hệ số tương quan tuyến tính giữa các tài sản. Công thức này cho thấy rủi ro danh mục không phụ thuộc vào biến thiên riêng lẻ của từng tài sản mà phụ thuộc vào mối tương quan giữa các tài sản đó. Tuy nhiên, $\rho$ chỉ đo lường mối quan hệ tuyến tính, và đây chính là điểm yếu mà hàm Copula sẽ khắc phục bằng cách tách biệt hàm phân phối biên và mô phỏng cấu trúc phụ thuộc phi tuyến.

Khi xem xét hai trường hợp biên với giả sử một danh mục gồm $N$ tài sản có trọng số bằng nhau $w_i = \frac{1}{N}$, ta có các kết quả cụ thể. Trong trường hợp thứ nhất khi các tài sản không tương quan và độc lập với nhau, tức là $\rho_{ij} = 0$ với mọi $i \neq j$, phương sai danh mục được tính bằng:
$$\text{Var}(R_w) = \sum_{i=1}^{N} \left(\frac{1}{N}\right)^2 \sigma_i^2 = \frac{1}{N^2} \sum_{i=1}^{N} \sigma_i^2$$
Khi $N \to \infty$ thì $\text{Var}(R_w) \to 0$. Việc có nhiều tài sản độc lập dẫn đến rủi ro danh mục càng nhỏ, điều này cho thấy rủi ro phi hệ thống có thể được triệt tiêu hoàn toàn nếu các tài sản không có sự liên kết với nhau. Trong trường hợp thứ hai khi các tài sản có tương quan với nhau, tức là $\rho_{ij} = c > 0$ với mọi $i \neq j$, hiệp phương sai giữa mọi cặp tài sản là $\sigma_{ij} = \text{Cov}(R_i, R_j) = c\sigma^2$. Khi đó, phương sai danh mục trở thành:
$$\begin{aligned}
\text{Var}(R_w) &= \frac{1}{N^2} \sum_{i=1}^{N} \sum_{j=1}^{N} \sigma_{ij} = \frac{1}{N^2} \left( \sum_{1 \le i \le N} \sigma_{ij} + \sum_{1 \le i < j \le N} \sigma_{ij} \right) \\
&= \frac{\sigma^2}{N} + \left(1 - \frac{1}{N}\right)c\sigma^2 = \frac{1-c}{N}\sigma^2 + c\sigma^2
\end{aligned}$$
Khi $N \to \infty$, phương sai danh mục sẽ hội tụ về giới hạn:
$$\lim_{N \to \infty} \text{Var}(R_w) = c\sigma^2$$
Như vậy, dù có tăng số lượng tài sản đến vô hạn thì phương sai danh mục vẫn không thể nhỏ hơn $c\sigma^2$. Nếu $c = 0$ đối với tài sản độc lập, kết quả sẽ trở về trường hợp thứ nhất với rủi ro giảm dần về $0$ khi $N$ tăng. Ngược lại, khi $c > 0$, rủi ro hệ thống chung giữa các tài sản luôn tồn tại và không thể loại bỏ được bằng đa dạng hóa. Hạn chế của mô hình Markowitz là rủi ro này phụ thuộc chặt chẽ vào hệ số $c$ tương quan tuyến tính. Tuy nhiên, trên thực tế, hệ số $c$ thường tăng mạnh trong các giai đoạn thị trường hoảng loạn tạo nên hiện tượng tương quan đuôi. Nếu chỉ sử dụng mô hình trung bình - phương sai, nhà đầu tư sẽ đánh giá thấp rủi ro hệ thống này vì giả định $c$ là một hằng số tuyến tính không đổi.